\documentclass[conference]{IEEEtran}
\usepackage{times}
\usepackage{amsmath,amssymb}
\usepackage{graphicx}
\usepackage{booktabs}
\usepackage{multirow}
\usepackage{url}
\usepackage{algorithm}
\usepackage{algpseudocode}
\usepackage{xcolor}

\title{Autonomous Self-Learning Serverless Intelligent Firewall:\\
Integrating REST API-Driven Open-Source Threat Intelligence,\\
Multi-Paradigm Machine Learning, and Federated Zero-Trust Architectures}

\author{\IEEEauthorblockN{Md Anisur Rahman Chowdhury and Kefei Wang}
\IEEEauthorblockA{Department of Computer and Information Science, Gannon University, USA\\
\texttt{engr.aanis@gmail.com}}}

\begin{document}
\maketitle

\begin{abstract}
This paper presents ASLF-OSINT, a serverless intelligent firewall architecture that unifies live open-source threat intelligence, multi-paradigm autonomous learning, and federated zero-trust policy control. The design introduces a super control plane for global governance and compact tenant-local subsystems for resilient local enforcement. Reported results include 98.7\% detection accuracy, 87 ms mean response latency, 94.3\% few-shot zero-day detection, 12.4-minute drift adaptation, and 99.8\% federated policy consistency.
\end{abstract}

\begin{IEEEkeywords}
serverless security, autonomous learning, threat intelligence, federated learning, zero trust, multi-tenant security
\end{IEEEkeywords}

\section{Introduction}
Serverless systems require adaptive security beyond static signatures and periodic retraining. Workloads are ephemeral, event-driven, and distributed across heterogeneous trust domains, making centralized-only controls fragile. Earlier research versions established baseline intrusion detection and cross-cloud orchestration, but they still depended on static data snapshots and centralized update assumptions. ASLF-OSINT extends this lineage with continuous threat-intelligence ingestion, autonomous adaptation, and multi-tenant operational governance.

\section{Problem Statement and Objectives}
Modern organizations require a security platform that can:\
(1) continuously ingest emerging threat indicators, \
(2) adapt to distribution shifts without full manual retraining, \
(3) support many enterprise tenants under one operator, and \
(4) maintain local enforcement even under temporary control-plane instability.

The objective is to deliver a practical, reproducible architecture where a provider-operated super control system manages tenant-specific compact subsystems with synchronized policy and version evolution.

\section{System Architecture}
\subsection{Super Control Plane}
The super control layer performs five functions: tenant account lifecycle management, policy publication, model/version rollout, federated telemetry aggregation, and global posture analytics.

\subsection{Tenant Compact Subsystem}
Each tenant receives a compact local protection subsystem that enforces allow/challenge/block actions over local networks and cloud services. The subsystem maintains enforcement continuity when central connectivity is degraded and re-synchronizes policy/model updates upon reconnect.

\subsection{Threat Intelligence Ingestion}
The ingestion layer queries MISP, AlienVault OTX, and VirusTotal through REST endpoints, then normalizes indicators (IP, domain, hash) into unified feature context for runtime scoring and scheduled retraining.

\subsection{Federated Zero-Trust Flow}
The decision path combines risk scoring with identity context and policy rules, then maps actions to provider-native controls across AWS, Azure, GCP, and private nodes.

\section{Learning Methodology}
ASLF-OSINT integrates five learning paradigms:
\begin{itemize}
  \item \textbf{Reinforcement learning (PPO):} improves action policy under latency and false-positive constraints.
  \item \textbf{Continual adaptation (SSF):} tracks and responds to concept drift.
  \item \textbf{Few-shot meta-learning:} improves zero-day recognition from limited labels.
  \item \textbf{Transfer learning:} reuses representations across heterogeneous domains.
  \item \textbf{Federated aggregation (FedNova):} updates distributed models while preserving tenant data locality.
\end{itemize}

\subsection{Risk Formulation}
For event $e_t$, the overall risk score is:
\begin{equation}
R_t = w_b B_t + w_a A_t + w_i I_t + w_o O_t + \Delta_t
\end{equation}
where $B_t$ is behavior risk, $A_t$ anomaly risk, $I_t$ identity risk, $O_t$ external intelligence risk, and $\Delta_t$ is severe-event adjustment. Weights are updated through feedback-driven adaptation.

\section{Implementation and Deployment}
\subsection{Execution Stack}
The implementation uses Python-based services with standard-library HTTP interfaces for portability, tenant-aware API endpoints, and browser dashboards for super and tenant operations.

\subsection{Core Endpoints}
\begin{itemize}
  \item Super control: tenant provisioning, asset onboarding, upgrade publication, fleet dashboard.
  \item Tenant subsystem: event protection, sync operation, tenant dashboard.
  \item Authentication: JWT issuance for super and tenant roles.
  \item Streaming: WebSocket channels for real-time super/tenant dashboard updates.
\end{itemize}

\subsection{Orchestration Algorithm}
\begin{algorithm}[t]
\caption{Tenant-Local Autonomous Protection with Super Sync}
\begin{algorithmic}[1]
\Procedure{ProtectEvent}{$e$}
  \State $o \gets \textsc{QueryOSINT}(e)$
  \State $x \gets \textsc{BuildFeatures}(e, o)$
  \State $r \gets \textsc{PredictRisk}(x)$
  \State $d \gets \textsc{ZeroTrustDecision}(r, e.identity, e.policy)$
  \State \textsc{ApplyProviderAction}$(d)$
  \State \textsc{ForwardTelemetryToSuper}(e, r, d)
  \State \Return $d$
\EndProcedure
\end{algorithmic}
\end{algorithm}

\section{Experimental Setup}
\textbf{Datasets:} CICIDS2017, CIC-IDS2018, UNSW-NB15, NSL-KDD.\\
\textbf{Live feeds:} 30 days of MISP, OTX, and VirusTotal indicators.\\
\textbf{Runtime targets:} AWS Lambda, Azure Functions, Google Cloud Functions, and Proxmox federation nodes.\\
\textbf{Evaluation axes:} detection quality, zero-day recognition, drift adaptation delay, latency, policy consistency, and analyst trust from explanation outputs.

\section{Results}
\begin{table}[t]
\centering
\caption{Key Reported Performance}
\begin{tabular}{ll}
\toprule
Metric & Result \\
\midrule
Detection accuracy & 98.7\% \\
Average response latency & 87 ms \\
Few-shot zero-day detection & 94.3\% \\
Concept drift adaptation time & 12.4 min \\
Federated policy consistency & 99.8\% \\
Analyst trust score (explainability) & 91.2\% \\
\bottomrule
\end{tabular}
\end{table}

\section{Ablation Study}
\begin{table}[t]
\centering
\caption{Ablation of Architecture Components}
\begin{tabular}{lcc}
\toprule
Configuration & Accuracy & Zero-day \\
\midrule
Baseline LSTM path only & 96.1\% & 71.2\% \\
+ Cross-cloud orchestration & 97.3\% & 78.6\% \\
+ OSINT ingestion & 98.0\% & 86.9\% \\
+ Continual/meta/federated stack (full) & 98.7\% & 94.3\% \\
\bottomrule
\end{tabular}
\end{table}

The ablation indicates that live OSINT context and autonomous adaptation jointly contribute most strongly to zero-day and drift-sensitive performance.

\section{Threat Model and Security Assumptions}
\subsection{Adversary Capabilities}
The adversary can generate high-volume traffic, attempt credential abuse, exploit policy drift windows, and probe cross-cloud inconsistency.

\subsection{Assumptions}
\begin{itemize}
  \item Tenant subsystems maintain trusted runtime integrity.
  \item Control-plane channels are TLS protected.
  \item Policy and version bundles are authenticated before application.
\end{itemize}

\subsection{Residual Risks}
Potential residual risks include compromised tenant credentials, poisoned threat-intelligence sources, and prolonged partition scenarios. These are mitigated through token rotation, source confidence weighting, and local enforcement fallback.

\section{Ethics, Compliance, and Privacy}
The architecture is designed for defensive monitoring and policy enforcement only. It minimizes tenant data movement via federated update patterns and avoids exposing raw tenant traffic to other tenants. Access to protected research artifacts is policy-gated, and system logs should be handled under applicable privacy and organizational compliance policies.

\section{Reproducibility Appendix}
\subsection{Artifact Layout}
\begin{itemize}
  \item Core implementation: \texttt{src/sif/*}
  \item Multi-tenant logic: \texttt{src/sif/multi\_tenant.py}
  \item Secure API and streaming: \texttt{src/sif/api\_server.py}
  \item Dashboards: \texttt{docs/super-dashboard.html}, \texttt{docs/tenant-dashboard.html}
  \item Tests: \texttt{tests/test\_*.py}
\end{itemize}

\subsection{Minimal Reproduction Steps}
\begin{enumerate}
  \item Start API server with configured JWT secret and admin credentials.
  \item Authenticate via \texttt{/auth/super/login} and create a tenant.
  \item Authenticate tenant via \texttt{/auth/tenant/login}.
  \item Submit events to \texttt{/tenant/\{tenant\_id\}/events}.
  \item Observe real-time updates through \texttt{/ws/super} and \texttt{/ws/tenant/\{tenant\_id\}}.
\end{enumerate}

\section{Conclusion}
ASLF-OSINT advances serverless firewall research by integrating autonomous adaptation, federated governance, and enterprise-ready multi-tenant operation. The framework demonstrates how super-to-tenant architectural separation can improve operational resilience while preserving centrally coordinated policy evolution.

\bibliographystyle{IEEEtran}
\bibliography{references}
\end{document}
